\subsection{The four-unitary free energy: a controlled local calculation}
\label{dev:1982-migdal}

\paragraph{The question actually recorded.}
Item \hyperref[item:1982-565.6]{1982-565.6}, problem
\hyperref[prob:GP-1982-02]{GP-1982-02}, is introduced on
\GSpage{P001040}, equation~(19), by \ru{Надо вычислить}.
Both the Russian transcription and an independent inspection of the scan show
a negative sign in the exponential, an ordered sum over \(\mu,\nu\), and an
inserted \(N^{-2}\log\). The notes ask for the limit; they do not supply a
proof of its existence, a domain for the coupling, or a normalization of Haar
measure. We make the last choice explicit by taking probability Haar measure.
The strongest formulation directly supported by the page is: evaluate the
displayed large-matrix limit on its domain of existence. No continuum-limit
claim follows simply from writing this integral.

For fixed \(d\geq2\), independent Haar matrices \(U_1,\ldots,U_d\in U(N)\),
and real \(\beta\), put
\begin{align}
 S_{d,N}(U)&=\sum_{\mu,\nu=1}^{d}
       \operatorname{Tr}(U_\mu U_\nu U_\mu^*U_\nu^*),\label{eq:1982-S}\\
 F_{d,N}(\beta)&=\frac1{N^2}\log
       \mathbb E\exp\{-\beta N S_{d,N}\}.\label{eq:1982-F}
\end{align}
The notebook case is \(d=4\). The general \(d\) notation is our extension
for checking coefficients. Pairing \((\mu,\nu)\) with \((\nu,\mu)\) makes
\(S_{d,N}\) real. The diagonal terms give \(dN\), exactly. Thus removing
them adds \(d\beta\) to \(F_{d,N}\); reversing the action sign replaces
\(\beta\) by \(-\beta\). These changes cannot be suppressed when comparing
formulas. In particular, the usual positive plaquette coupling corresponds
to negative \(\beta\) in \eqref{eq:1982-F}.

\paragraph{What later work establishes.}
The reduced lattice model itself predates this seminar:
\citet{eguchikawai1982reduction} was published on 19~April~1982.
That identification does not evaluate the four-matrix integral.
For our exact mathematical normalization,
\citet[Theorem~0.1]{collinsguionnetmaurel2009asymptotics}
is decisive: for a fixed self-adjoint noncommutative polynomial potential in
finitely many independent Haar unitaries, the normalized logarithmic integral
has a large-\(N\) limit for sufficiently small real coupling; that limit is
analytic near zero and its Taylor coefficients are the limits of the formal
finite-\(N\) coefficients. There are no external matrices here, so the
boundedness and joint-moment assumptions on those matrices are vacuous.
Take the potential to be the negative of the ordered commutator sum in
\eqref{eq:1982-S}. The preprint first circulated on 8~August~2006; publication
followed in 2009.

The distinction between an unmodified reduced model and a valid reduction
of an extended gauge theory remains material. Center symmetry and the precise
quenching prescription are discussed in
\citet[Secs.~II--III]{neuberger2020quenched}; neither quenching nor twisting
is present in the notebook's integral. The mathematical local convergence
theorem therefore applies directly to the displayed integral, while a
physical equivalence requires separate hypotheses. The search behind this
chapter located no verified all-coupling evaluation of this exact four-unitary
integral. This is a bounded literature finding as of November~2026, not a
proof that no such evaluation exists. The sources and failed leads are
recorded in the research audit.
The checked sources do not establish that the exact integral was abandoned.
They do explain why its evaluation and its use as a reduction of QCD are
different research tasks: the latter depends on symmetry and on the chosen
reduction prescription, while the former still involves four interacting
unitary matrices.

\paragraph{The stronger target and the scope of this development.}
A plausible stronger formulation is existence and determination of
\(F_4(\beta)=\lim_{N\to\infty}F_{4,N}(\beta)\) for every real \(\beta\),
including any points where it fails to be differentiable. This is our
formulation of the global target; it is not a theorem asserted by the notes
or established in this chapter. We supply exact finite-\(N\) coefficients,
elementary global bounds, and the rigorously justified local expansion.
These are reproducible bounded deductions, not claims of a new solution
of the model or of priority for the coefficient formulas.

\begin{proposition}[Normalization, convexity, and compactness \ours]
For every \(d,N\) and every real \(\beta\), \(F_{d,N}\) is well-defined,
real analytic, convex, and \(d^2\)-Lipschitz. It satisfies
\[
 -d\beta-d(d-1)|\beta|\ \leq\ F_{d,N}(\beta)
 \ \leq\ -d\beta+d(d-1)|\beta|.
\]
Every sequence of matrix sizes tending to infinity has a subsequence along
which \(F_{d,N}\) converges uniformly on every compact real interval to a
convex \(d^2\)-Lipschitz function vanishing at zero.
\end{proposition}
\begin{proof}
Each commutator is unitary, hence its trace has absolute value at most \(N\).
The off-diagonal sum is real and has absolute value at most \(d(d-1)N\).
Applying these pointwise bounds inside the positive integral proves the
inequalities. At any real coupling the integrand is bounded on a compact
group, so differentiation under the integral is legitimate to every order.
The partition function is entire in complex coupling and strictly positive
on the real axis, hence its logarithm is analytic in a neighborhood of each
real point. It need not have an entire logarithm. With expectation under
the tilted probability measure denoted by \(\mathbb E_\beta\),
\[
 F_{d,N}'(\beta)=-\frac1N\mathbb E_\beta S_{d,N},\qquad
 F_{d,N}''(\beta)=\operatorname{Var}_\beta(S_{d,N})\geq0.
\]
The bound \(|S_{d,N}|\leq d^2N\) gives the Lipschitz constant.
Arzel\`a--Ascoli on each interval \([-k,k]\), followed by a diagonal
subsequence, gives compact convergence. Convexity and the Lipschitz bound
pass to that limit. This argument gives subsequences only, not uniqueness
of the limiting free energy.
\end{proof}

\begin{lemma}[One plaquette and two adjacent plaquettes \ours]
Let \(U,V,W\) be independent probability-Haar matrices in \(U(N)\),
\(N\geq2\). Write \(C(U,V)=\operatorname{Tr}(UVU^*V^*)\) and
\(X(U,V)=C(U,V)+\overline{C(U,V)}\). Then
\begin{align*}
 \mathbb E C(U,V)&=N^{-1},&
 \mathbb E C(U,V)^2&=\frac4{N^2-1},&
 \mathbb E|C(U,V)|^2&=\frac{N^2}{N^2-1},\\
 \operatorname{Var}X(U,V)&=
       \frac{2(N^2+4)}{N^2-1}-\frac4{N^2},&
 \operatorname{Cov}(X(U,V),X(U,W))&=\frac4{N^2}.
\end{align*}
Plaquettes involving disjoint pairs of matrices are independent.
\end{lemma}
\begin{proof}
The first Haar identity is
\(\mathbb E_U UAU^*=N^{-1}\operatorname{Tr}(A)I\).
For fixed matrices \(A,B,C,D\), the second identity gives
\begin{align*}
 &\mathbb E_U\bigl[\operatorname{Tr}(UAU^*B)
                    \operatorname{Tr}(UCU^*D)\bigr]\\
 &\quad=\frac{\operatorname{Tr}A\operatorname{Tr}B
                    \operatorname{Tr}C\operatorname{Tr}D
                    +\operatorname{Tr}(AC)\operatorname{Tr}(BD)}{N^2-1}\\
 &\qquad\ -\frac{\operatorname{Tr}(AC)\operatorname{Tr}B\operatorname{Tr}D
                    +\operatorname{Tr}A\operatorname{Tr}C\operatorname{Tr}(BD)}
                    {N(N^2-1)}.
\end{align*}
This is the degree-two specialization of
\citet[Corollary~2.4]{collinssniady2006integration}; it also follows by
projecting \(A\otimes C\) onto the span of the identity and the tensor flip,
then solving the two equations obtained by taking its trace and flip-trace.
The determinant of those two equations is \(N^2(N^2-1)\), explaining
the restriction \(N\geq2\).

Put \(t=\operatorname{Tr}V\) and \(s=\operatorname{Tr}(V^2)\).
Schur orthogonality, using
\(t^2=\chi_{(2)}+\chi_{(1,1)}\) and
\(s=\chi_{(2)}-\chi_{(1,1)}\), yields
\[
 \mathbb E|t|^2=1,\quad \mathbb E|t|^4=2,\quad
 \mathbb E|s|^2=2,\quad \mathbb E(s\overline t^{\,2})=0.
\]
Both degree-two characters exist and are distinct for \(N\geq2\).
First condition on \(V\). Substituting \(A=C=V\), \(B=D=V^*\)
into the second identity gives
\[
 \mathbb E_U C(U,V)^2
 =\frac{|t|^4+|s|^2}{N^2-1}
   -\frac{s\overline t^{\,2}+\overline s t^2}{N(N^2-1)}.
\]
For the absolute square, use
\(\overline{C(U,V)}=\operatorname{Tr}(UV^*U^*V)\); the result is
\[
 \mathbb E_U|C(U,V)|^2
     =\frac{|t|^4+N^2-2|t|^2}{N^2-1}.
\]
Averaging the displayed formulas proves the three moments and hence the
variance. Finally, conditioning on the shared matrix \(U\), the two
plaquettes are independent and each has conditional mean
\(2|\operatorname{Tr}U|^2/N\). Their covariance is consequently
\(4\operatorname{Var}(|\operatorname{Tr}U|^2)/N^2=4/N^2\).
Independence of disjoint pairs is immediate from product Haar measure.
\end{proof}

\begin{proposition}[Exact first and second derivatives \ours]
For \(d\geq2\) and \(N\geq2\),
\begin{align}
 F_{d,N}'(0)&=-d-\frac{d(d-1)}{N^2},\label{eq:1982-first}\\
 F_{d,N}''(0)&=d(d-1)\left\{
      \frac{N^2+4}{N^2-1}+\frac{4d-10}{N^2}\right\}.
       \label{eq:1982-second}
\end{align}
For \(N=1\), the exact answer instead is
\(F_{d,1}(\beta)=-d^2\beta\).
\end{proposition}
\begin{proof}
The mean of \(S_{d,N}\) is
\(dN+d(d-1)/N\), which gives \eqref{eq:1982-first}.
There are \(m=d(d-1)/2\) unordered plaquettes and
\(a=d\binom{d-1}{2}\) unordered pairs of distinct plaquettes sharing a matrix.
The constant diagonal terms contribute no variance. The preceding lemma and
the factor of two for covariance terms in a variance give
\begin{align*}
 \operatorname{Var}S_{d,N}
 &=m\left\{\frac{2(N^2+4)}{N^2-1}-\frac4{N^2}\right\}
       +2a\frac4{N^2}\\
 &=d(d-1)\left\{\frac{N^2+4}{N^2-1}+\frac{4d-10}{N^2}\right\}.
\end{align*}
Since \(F_{d,N}''(0)=\operatorname{Var}S_{d,N}\), the formula follows.
At \(N=1\) every scalar commutator is one, so the integral is
\(\exp(-\beta d^2)\). No substitution into the singular \(N\geq2\)
formula is permissible.
\end{proof}

\begin{corollary}[The justified local answer for the notebook convention \ours]
There is an \(\varepsilon>0\) such that the large-\(N\) free energy in
the notebook exists for real \(|\beta|<\varepsilon\), has an analytic
extension near zero, and
\[
                 F_4(\beta)=-4\beta+6\beta^2+O(\beta^3).
\]
More generally \(F_d(\beta)=-d\beta+\tfrac12d(d-1)\beta^2+O(\beta^3)\)
locally. The error here is a statement about the analytic limiting function,
not a uniform finite-\(N\) error inferred from a finite computation.
\end{corollary}
\begin{proof}
Apply the local unitary-integral theorem stated above, then take the limits
of \eqref{eq:1982-first} and \eqref{eq:1982-second}.
The theorem, rather than the numerical or symbolic checks, justifies the
identification of these limits with derivatives of \(F_d\).
\end{proof}

\begin{remark}[Verification \oursComputed]
The independent calculation in \nolinkurl{code/haar_check.py} expands products of traces
into matrix indices and applies the degree-one and degree-two Weingarten
formulas directly. It counts free index loops with a union--find algorithm.
This does not use the conditional-covariance argument of the proof. Six
basic mixed moments agree as rational functions of the symbolic integer
\(N\geq2\); summing all 584 ordered plaquette pairs for
\(d=2,3,4,5\) independently reproduces \eqref{eq:1982-second}.
Additional exact rational evaluations use \(N=2,3,4,7,11\).
For the four-matrix model the second derivatives at \(N=2,3,4\) are
\(50,55/2,41/2\), respectively, and tend to \(12\).
The separate scalar check gives zero second derivative at \(N=1\).
The check confirms the finite-rank restriction and both coefficients;
it makes no inference from finite samples to a global analytic limit.
An independent review also rederived the single-plaquette moments from
\(\mathbb E\chi_\rho([U,V])=1/\dim\rho\): Schur averaging first in
\(U\) gives \(|\chi_\rho(V)|^2/\dim\rho\), and averaging in \(V\)
gives the identity. Applying it to the symmetric and alternating squares,
and to the decomposition of the fundamental representation tensored with
its dual, reproduces the same second moments for every \(N\geq2\).
That review separately checked the adjacent-pair count and the exact scope
of the cited local-limit theorem; it found no mathematical correction.
The source, symbolic output, and machine-readable rational results are
included in the archive.
\end{remark}

\paragraph{What remains and what another run could do.}
Neither the two coefficients nor subsequential compactness select a global
large-\(N\) branch. The remaining obstacle is control of the interacting
unitary ensemble outside the small-potential region, including competing
states and any transition points. A complete physical reduction is an
additional question. One further bounded run could calculate higher connected
Haar cumulants, provide explicit certified bounds on a small-coupling interval,
and compare conventions against the original lattice strong-coupling series.
The precise next calculation is the third and fourth Taylor coefficients of
the same ordered-sum integral with their finite-\(N\) corrections, followed by
a justified limiting argument. It would enlarge the controlled local answer;
there is no basis here for promising an all-coupling solution.
